\clearpage
\VerdierPhysicalPageAnchor{0175}%
\begingroup\small\setlength{\parskip}{0pt}\setlength{\emergencystretch}{1.5em}
\setlength{\abovedisplayskip}{4pt plus 2pt minus 2pt}%
\setlength{\belowdisplayskip}{4pt plus 2pt minus 2pt}%

\VerdierSemanticAnchor{ch3d-1-1-2-cont}
It follows immediately from the definition of distinguished triangles in
$K(\mathcal A)$ (Chap. I, 3.3.1) and the properties of the canonical
cohomological functor $H\colon K(\mathcal A)\longrightarrow\mathcal A$ that
the categories $K^{*,*'}(\mathcal A)\subset K(\mathcal A)$ are full
triangulated subcategories (Chap. II, 2.1.5), strictly full if $*={}$
``empty'', and saturated subcategories (Chap. II, 2.1.6) of
$K(\mathcal A)$.

\VerdierSemanticAnchor{ch3d-1-1-3}
\medskip\noindent\textbf{1.1.3.}\quad
A complex $X$ of $\mathcal A$ is said to be \emph{acyclic} if
$H^i(X)=0$ for every integer $i$. We denote by
$\operatorname{Ac}^{*}(\mathcal A)$ the full subcategory of
$K^{*}(\mathcal A)$ defined by the acyclic complexes that are objects of
$K^{*}(\mathcal A)$. A morphism $X\xrightarrow{u}Y$ of $K(\mathcal A)$
(respectively, of $\operatorname{comp}(\mathcal A)$) is called a
\emph{quasi-isomorphism} if, for every integer $i$, the morphism
$H^i(X)\xrightarrow{H^i(u)}H^i(Y)$ is an isomorphism. We denote by
$\operatorname{Qis}^{*}(\mathcal A)$ (or by $\operatorname{Qis}^{*}$ when
no confusion can result) the set of quasi-isomorphisms of
$K^{*}(\mathcal A)$. The subcategory $\operatorname{Ac}^{*}(\mathcal A)$
of $K^{*}(\mathcal A)$ is a strictly full, saturated triangulated
subcategory of $K^{*}(\mathcal A)$. The set
$\operatorname{Qis}^{*}(\mathcal A)$ is a saturated multiplicative system
of morphisms of $K^{*}(\mathcal A)$ (Chap. II, 2.1.4). In a distinguished
triangle $(X,Y,Z,u,v,w)$ of $K^{*}(\mathcal A)$, the morphism $u$ is a
quasi-isomorphism if and only if the complex $Z$ is acyclic. In other words
(Chap. II, 2.1.15), $\operatorname{Qis}^{*}(\mathcal A)$ is the
multiplicative system of $K^{*}(\mathcal A)$ associated with the full
triangulated subcategory $\operatorname{Ac}^{*}(\mathcal A)$.

\VerdierSemanticAnchor{ch3d-1-1-4}
\medskip\noindent\textbf{1.1.4.}\quad
Let $M$ be a set of objects of $\mathcal A$. A complex of $\mathcal A$ is
said to be \emph{of type $M$} if all its components are elements of $M$. A
morphism $u\colon X\longrightarrow Y$ of $K^{*}(\mathcal A)$ or of
$\operatorname{comp}^{*}(\mathcal A)$ is called a \emph{right resolution}
(respectively, a \emph{left resolution}) of $X$ (respectively, $Y$)
\emph{of type $M$} if $u$ is a quasi-isomorphism and if $Y$ (respectively,
$X$) is a complex of type $M$. Most often, by abuse of language and when
the context leaves no room for confusion, the qualification ``right''
(respectively, ``left'') is omitted. Likewise, most often, by abuse of
language and when the context leaves no room for confusion, the target
(respectively, the source) of a right (respectively, left) resolution of an
object $Y$ is itself called a right (respectively, left) resolution of
$Y$. We denote by $\operatorname{comp}_{\mathcal A}^{*}(M)$, or simply
$\operatorname{comp}^{*}(M)$ (respectively, $K_{\mathcal A}^{*}(M)$, or
simply $K^{*}(M)$), the full subcategory of
$\operatorname{comp}^{*}(\mathcal A)$ (respectively, $K^{*}(\mathcal A)$)
defined by the complexes of type $M$.

\VerdierSemanticAnchor{ch3d-1-1-5}
\medskip\noindent\textbf{1.1.5.}\quad
When $M$ is the set of projective (respectively, injective) objects of the
category $\mathcal A$, complexes whose components are projective
(respectively, injective) are called, in accordance with the conventions
above, \emph{complexes of projective type} (respectively, \emph{injective
type}). The reader should take care not to confuse complexes of projective
type (respectively, injective type) with projective (respectively,
injective) objects of the category $\operatorname{comp}(\mathcal A)$. One
checks easily that an object of $\operatorname{comp}(\mathcal A)$ is
projective (respectively, injective) if and only if it is of projective type
(respectively, injective type) and homotopic to zero.
\endgroup
